\section{问答精读：从课堂观点到团队 SOP}

除了主讲内容，问答环节给了很多工程化细节，尤其是如何把 ``Virtual Lab + Paper2Agent + MCP'' 真正落地成可持续流程。本章把这些问答信息整理成可执行 SOP，避免只停留在概念层面。

\subsection{如何触发 Agent 之间的自主协作}

针对``Agent-to-Agent 对话如何触发''的问题，课程回答是通过 paper MCP 的标准接口，把多个论文能力挂到同一协调器（通常是 chatbot 或 orchestration agent）上，再由协调器识别任务中可组合的资源并发起调用。

这与很多团队的常见做法不同。常见做法是让单个 Agent 靠长上下文``知道所有东西''，最终导致上下文拥挤、路由不清晰、可解释性差。课程给出的做法是能力模块化与调用协议化，使协作触发条件可以显式配置。

\begin{knowledgebox}{最小可用协作触发器（建议实现）}
\begin{itemize}
    \item 定义每个 paper MCP 的能力描述：输入类型、输出类型、前置条件。
    \item 在协调器中实现``任务需求 $\rightarrow$ 能力匹配''规则。
    \item 允许协作链路有拒绝分支：能力不满足时直接回退给人类。
    \item 对每次跨 MCP 调用记录 trace，用于后续失效分析。
\end{itemize}
\end{knowledgebox}

\subsection{关于上下文成本与 token 开销}

问答还提到 MCP 的 token overhead 相对较低，核心原因是很多上下文被移到了工具与资源层，而不是全部塞进 prompt。对工程团队来说，这意味着优化重点应放在\textbf{接口设计}与\textbf{资源检索准确率}，而不是单纯压缩提示词。

\begin{importantbox}{控制上下文成本的三条策略}
\begin{itemize}
    \item 把可执行能力放进 MCP 工具，不把流程细节重复写进会话上下文。
    \item 先检索再拼接：只把当前任务需要的 paper 片段拉入上下文。
    \item 将中间状态结构化存储，避免每轮对话重复传输历史长文本。
\end{itemize}
\end{importantbox}

如果团队只盯着``本轮 token 少了多少''，而忽略了失败重试次数与人工兜底成本，通常会做出局部最优决策。课程案例里 Paper2Agent 的优势之一正是减少失败回滚，降低总体执行时间。

\subsection{证据时间轴：关键结论对应到视频片段}

为方便后续复核与组内复盘，表\ref{tab:evidence-index} 把本讲关键论点映射到时间戳、证据类型和可执行动作。这样团队在二次讨论时可以直接回看对应片段，而不是凭记忆复述。

\begin{table}[H]
\centering
\small
\begin{tabular}{p{0.13\textwidth}p{0.25\textwidth}p{0.19\textwidth}p{0.31\textwidth}}
\toprule
时间戳 & 课堂信息点 & 证据类型 & 可执行动作 \\
\midrule
00:01--00:03 & AI 从工具转向 co-scientist & 主讲开场定义 & 统一团队术语，区分工具层与协作层目标 \\
00:03--00:06 & Virtual Lab 角色与会议机制 & 系统结构说明 & 为本地项目定义 PI/Student/Critic 最小角色集 \\
00:05:32 & Agent 创建 + 组会 + 单聊流程图 & 幻灯片关键帧 & 按图实现路由器、并行会话、汇总器三模块 \\
00:06--00:09 & nanobody 路线与 Critic 介入 & 案例讨论文本 & 在高风险任务默认开启反方 Agent \\
00:14:40 & Agent School 学习流程 & 幻灯片关键帧 & 建立周度主题学习队列与能力回归测试 \\
00:18:54 & 候选结果实验验证示意 & 幻灯片关键帧 & 为每轮候选绑定实验优先级与验证预算 \\
00:21--00:25 & memory/sandbox/HITL 设计 & 问答与方法总结 & 对所有任务持久化记忆与可重放执行日志 \\
00:26--00:30 & Paper2Agent 目标与流程 & 主讲流程讲解 & 选择高价值论文先做 MCP 化 \\
00:29:27 & 环境/抽取/测试 Agent 管线 & 幻灯片关键帧 & 先搭建三 Agent 流水线再做下游 UI \\
00:32--00:33 & 评测中相对 baseline 更稳更快 & benchmark 口述结果 & 用同任务 A/B 测试验证本地收益 \\
00:34--00:37 & Agent-to-Agent 协作发现新线索 & 协作案例讲解 & 为方法 MCP 与数据 MCP 建自动匹配器 \\
00:47--00:48 & 300+ 投稿、48 篇录用 & 会议统计口述 & 参考其 checklist 设计内部项目评审模板 \\
00:50--00:52 & 多模型评审风格差异 & 评审分布观察 & 使用多评审 Agent + 校准，不用单模型定生死 \\
00:52:30 & 评审可抓错但会 sycophancy & 幻灯片关键帧 & 给评审 Agent 加反谄媚提示与硬约束模板 \\
00:53--00:54 & 参考文献幻觉核验结果 & 自动检查统计 & 在稿件提交流程中前置引用核验 \\
00:58--01:00 & 局限：原创能力与上下文边界 & 主讲反思 & 把人类聚焦在方法创新和前提设定 \\
01:00--01:01 & 强调闭环人机协作 & 课程收束观点 & 固化``计算-实验-回流''闭环到团队 SOP \\
\bottomrule
\end{tabular}
\caption{课程证据索引：论点、时间戳与可执行动作映射}
\label{tab:evidence-index}
\end{table}

\subsection{团队落地模板：四周试点方案}

为了把本讲快速转成组织收益，可以采用一个四周试点：

\begin{itemize}
    \item 第 1 周：选题与基线。选择 1--2 个复现负担高的论文任务，记录人工 baseline 时间与错误类型。
    \item 第 2 周：Paper2Agent 化。完成 environment/extraction/testing 三 Agent 管线，产出首版 MCP。
    \item 第 3 周：接入 Virtual Lab。给 MCP 配置 PI、两个 specialist 和 Critic，跑两轮并行会议。
    \item 第 4 周：复盘与治理。对比成功率、耗时、人工介入点；补充引用核验和 human sign-off 规则。
\end{itemize}

\begin{warningbox}{试点常见失败模式}
\begin{itemize}
    \item 一上来就追求全自动，导致边界条件未定义清楚。
    \item 只记录最终答案，不记录中间调用与失败重试轨迹。
    \item 缺少统一评价指标，最后无法判断是否优于人工 baseline。
\end{itemize}
\end{warningbox}

建议在试点前把成功标准写死：例如任务完成率、平均耗时、人工审阅时长、引用错误数、可复现实验比例。这样在汇报阶段才能客观回答``这套系统值不值得继续投入''。

\subsection{本章小结}

问答环节的最大价值是把理念转成工程动作：协作触发要协议化、上下文成本要系统性看总账、评审与引用核验要前置到流程。按证据时间轴执行，团队可以在一个月内完成第一轮可衡量试点。

